\newpage
\phantomsection
\label{prf:monotone-continuity-criterion}

\begin{remark*}[Return]
\hyperref[cor:monotone-continuity-criterion]{Return to Corollary}
\end{remark*}

\begin{theorem*}[Continuity Criterion for Monotone Functions]
Let $f : I \to \mathbb{R}$ be monotone increasing and let $c \in I$ be an interior point. Then $f$ is continuous at $c$ if and only if $f(c^{-}) = f(c^{+})$.
\end{theorem*}

\begin{proof}
\textbf{Professional Standard Proof.}~\\
Let $f : I \to \mathbb{R}$ be monotone increasing and let $c \in I$ be an interior point.

$(\Rightarrow)$ Suppose $f$ is continuous at $c$. Then $\lim_{x \to c} f(x)$ exists and equals $f(c)$. By the criterion that a two-sided limit exists if and only if both one-sided limits exist and are equal, both $\lim_{x \to c^{-}} f(x)$ and $\lim_{x \to c^{+}} f(x)$ exist and equal $f(c)$. In particular $f(c^{-}) = f(c^{+})$.

$(\Leftarrow)$ Suppose $f(c^{-}) = f(c^{+})$. By the One-Sided Limits of Monotone Functions theorem, these one-sided limits equal $\sup\{f(x) : x < c\}$ and $\inf\{f(x) : x > c\}$ respectively, and the same theorem gives the ordering $f(c^{-}) \leq f(c) \leq f(c^{+})$. Since $f(c^{-}) = f(c^{+})$, this ordering forces $f(c^{-}) = f(c) = f(c^{+})$. Hence both one-sided limits exist and equal $f(c)$, so the two-sided limit $\lim_{x \to c} f(x) = f(c)$ exists, and $f$ is continuous at $c$.
\end{proof}

\begin{proof}
\textbf{Detailed Learning Proof.}~\\

\textbf{Step 1.} Establish the forward direction: continuity implies matching one-sided limits. If $f$ is continuous at $c$, then the two-sided limit $\lim_{x \to c} f(x)$ exists and equals $f(c)$. By the characterization of two-sided limits in terms of one-sided limits, this means both $\lim_{x \to c^{-}} f(x)$ and $\lim_{x \to c^{+}} f(x)$ exist and both equal $f(c)$. Therefore $f(c^{-}) = f(c) = f(c^{+})$, so in particular $f(c^{-}) = f(c^{+})$.

\textbf{Step 2.} Establish the reverse direction: matching one-sided limits imply continuity. Assume $f(c^{-}) = f(c^{+})$. By the One-Sided Limits of Monotone Functions theorem, both one-sided limits exist and satisfy the ordering $f(c^{-}) \leq f(c) \leq f(c^{+})$. Since $f(c^{-}) = f(c^{+})$ by hypothesis, this ordering forces all three values to be equal: $f(c^{-}) = f(c) = f(c^{+})$.

\textbf{Step 3.} Conclude continuity from matching one-sided limits. Since both one-sided limits exist and equal $f(c)$, the criterion for two-sided limits guarantees that $\lim_{x \to c} f(x)$ exists and equals $f(c)$. This is precisely the definition of continuity at $c$.
\end{proof}

\begin{remark*}[Proof structure]
A direct biconditional proof. The forward direction unpacks the definition of continuity via the equivalence of two-sided and matching one-sided limits. The reverse direction applies the ordering $f(c^{-}) \leq f(c) \leq f(c^{+})$ from the One-Sided Limits of Monotone Functions theorem: the hypothesis $f(c^{-}) = f(c^{+})$ squeezes $f(c)$ to the common value, after which the two-sided limit criterion closes the argument.
\end{remark*}

\begin{remark*}[Dependencies]~\\
\begin{itemize}
  \item \hyperref[def:continuous-at-point]{Definition of Continuity at a Point}
  \item \hyperref[thm:monotone-one-sided-limits]{One-Sided Limits of Monotone Functions}
  \item \hyperref[thm:two-sided-limit-iff-matching-one-sided-limits]{Two-Sided Limit Iff Matching One-Sided Limits}
\end{itemize}
\end{remark*}
